\subsubsection{Aviation}
\noindent
\begin{minipage}{.49\textwidth}
    \noindent\textbf{\LARGE Aviation\index{Aviation}}
\end{minipage}
\hspace{.02\textwidth}
\begin{minipage}{.49\textwidth}
    \raggedleft
    \textbf{\large Gin}\\
    \sweettartguage{0.4}\index[drinks]{Tart!40\%!Aviation}
\end{minipage}
\vspace{.5em}\hrule\vspace{.5em}
\noindent
\textbf{\large Flavor Profile:} Tart\index[flavor-profile]{Tart}, Dry\index[flavor-profile]{Dry}, Floral\index[flavor-profile]{Floral}, Cherry\index[flavor-profile]{Cherry}, Herbal\index[flavor-profile]{Herbal}, Light\index[flavor-profile]{Light}.
\vspace{.5em}\hrule\vspace{.5em}
\noindent
\textbf{\large History:} The Aviation was created by Hugo Ensslin, the head bartender at Hotel Wallick in New York in the early 1900s. It was considered one of the first "industry" drinks and if you ordered it, you were considered an insider in the cocktail world. The Aviation became less popular as crème de violette became more difficult to find and recipes often left it out. However, recently the original Aviation is making a comeback, and becoming a favorite once again.
\vspace{.5em}\hrule\vspace{.5em}
\noindent
\begin{minipage}[t]{.49\textwidth}
    \noindent
    \textbf{\large Ingredients}
    \begin{itemize}
        \item Gin 60ml\index[ingredient]{Gin}
        \item Maraschino 15ml\index[ingredient]{Maraschino}
        \item Lemon Juice 22,5ml\index[ingredient]{Lemon Juice}
        \item Crème de Violette 7,5ml\index[ingredient]{Crème de Violette}
    \end{itemize}
    \noindent
    \textbf{\large Garnish}
    \begin{itemize}
        \item Lemon Peel\index[ingredient]{Lemon Peel} or Dark Cherry\index[ingredient]{Dark Cherry}
    \end{itemize}
\end{minipage}
\hspace{.01\textwidth}\vrule\hspace{.01\textwidth}
\begin{minipage}[t]{.49\textwidth}
    \noindent
    \textbf{\large Tools/Glassware}
    \begin{itemize}
        \item Shaker
        \item Strainer & Fine Strainer (optional)
        \item Jigger
        \item Y-Peeler
        \item Coupe\index[glass]{Coupe} or Martini Glass\index[glass]{Martini Glass}
    \end{itemize}
\end{minipage}
\vspace{.5em}\hrule\vspace{.5em}
\noindent
\begin{minipage}{\textwidth}
    \noindent
    \textbf{\large Instructions}
    \begin{itemize}
        \item Place all ingredients into a shaker.
        \item Add ice to shaker and shake for 8--10 seconds.
        \item Double strain into coupe or martini glass.
        \item Garnish with lemon peel or cherry.
    \end{itemize}
\end{minipage}
\vspace{.5em}\hrule\vspace{.5em}
\noindent
\textbf{\large Fun Tip:} Crème de violette differs from brand to brand, and trying a new brand may slightly change the color of your cocktail.
\newpage
